
\textbf{Proposition 4.3.1.} \uline{Avec les notations ci-dessus, pour toute suite} \(\alpha_1,\ldots,\alpha_d\) \uline{d'entiers positifs avec} \(\sum_{i=1}^{d}\alpha_i=d\)\uline{, le coefficient de}
\[
 \binom{n_1}{\alpha_1}\cdots\binom{n_d}{\alpha_d}
\]
\uline{dans le polynôme de Snapper} \(\chi(L_1^{n_1}\cdots L_d^{n_d}F)\) \uline{est égal à}
\begin{equation}\tag{4.3.2}
 \bigl\langle c_1(L_1)^{\alpha_1}\cdots c_d(L_d)^{\alpha_d},\lvert\underline{F}\rvert\bigr\rangle,
\end{equation}
\uline{où} \(\lvert\underline{F}\rvert\) \uline{est la classe de} \(\underline{F}\) \uline{dans} \(\Gr_d(X)\)\uline{. Donc dans le développement ordinaire de ce polynôme suivant les monômes en des} \(n_i\) \((i=1,\ldots,d)\)\uline{, le coefficient du monôme} \(n_1^{\alpha_1}\cdots n_d^{\alpha_d}\) \uline{est égal à}
\[
 (1/(\alpha_1! \cdots a_d!))
 \bigl\langle c_1(L_1)^{\alpha_1}\cdots c_1(L_d)^{\alpha_d},\lvert\underline{F}\rvert\bigr\rangle.
\]

En particulier on obtient immédiatement le

\textbf{Corollaire 4.3.3.} \uline{Le coefficient du monôme} \(n_1n_2\cdots n_d\) \uline{dans le polynôme de Snapper} \(\chi(L_1^{\alpha_1}\cdots L_d^{\alpha_d}F)\) \uline{est égal à}
\begin{equation}\tag{4.3.4}
 \bigl\langle c_1(L_1)\cdots c_1(L_d),\lvert\underline{F}\rvert\bigr\rangle
 =
 (c_1(L_1)\cdots c_1(L_d),\lvert\underline{F}\rvert).
\end{equation}

\emph{Démonstration} : Posons \(L_i=1+M_i\) pour tout \(i=1,\ldots,d\). Tenant compte des définitions on trouve immédiatement que l'entier (4.3.2) est égal à \(\chi(M_1^{\alpha_1}\cdots M_d^{\alpha_d}F)\). Lorsqu'on substitue dans (2.1.4) \(A=K^{\bullet}(X)/\Filt^{d+1}(X)\), \(I=\Filt^1(X)/\Filt^{d+1}(X)\), alors (2.1.5) donne, (compte tenu de ce que \(M_i\in\Filt^1(X)\) pour tout \(i=1,\ldots,d\)) la congruence
\begin{equation}\tag{4.3.5}
 L_1^{n_1}\cdots L_d^{n_d}
 \equiv
 \sum_{\substack{\alpha_1,\ldots,\alpha_d\geq0\\ \alpha_1+\cdots+\alpha_d=d}}
 \binom{n_1}{\alpha_1}\cdots\binom{n_d}{\alpha_d}
 M_1^{\alpha_1}\cdots M_d^{\alpha_d}
 \quad\mathrm{mod.}\Filt^{d+1}(X).
\end{equation}
Donc le coefficient de \(\binom{n_1}{\alpha_1}\cdots\binom{n_d}{\alpha_d}\) dans le polynôme \(\chi(L_1^{n_1}\cdots L_d^{n_d}F)\) est égal à \(\chi(M_1^{\alpha_1}\cdots M_d^{\alpha_d}F)\), C.Q.F.D.

